% ============================================================
% M3 S2 导学件·波1 片件 —— 课时01 坐标法（2.1）＝全波母版（五臂照抄件）
% 生成：M3 成卷轮 S2 母版代理｜2026-09-14｜编译 cwd＝本目录（中文路径禁 \input 绝对路径）
%   xelatex main-true.tex ×2 → 含详解印本；xelatex main-false.tex ×2 → 纯题版（单源双本）
% 输入链（全只读）：题面库 课时01-坐标法.md（题面逐字装配）＋ -答案侧.md（值逐字回嵌）
%   ＋ 定稿/批A-课时01-坐标法.md（详解回嵌源）；toolchain＝qp-m3.sty（钉后版
%   md5 7c3930362be8a0a2bdf21bbf8ac16573，本地挂载副本，破损即停）。
% 母版体例：详见《母版体例说明.md》（结构骨架/宏用法/门谱跑法/常见坑）。
% 键账：件内 ansblock 键集 ≡ 题面库片键集 ≡ manifest 键序（23 键，零缺漏零多余）；
%   预习位零键零答案（口令④裁定前口径）；印面号 1—23＝探究 1—18＋课堂评价 19—23，
%   键↔号映射见 值台账-课时01.json。
% 括线判据（承重墙禁放宽）：估高＞8 行（chars/23 栏宽口径）→ 括线模，逐键读数入值台账；
%   其余灰底。本片括线键：E2、E3、E14、T1、T2（估高读数见台账）。
% 门谱读数：见 过程件 工作区/_tmpM3S2母版0914/（编译 log、门谱输出、目验 PNG 双档）。
% ============================================================
\documentclass[fontset=none]{ctexart}
\usepackage{qp-m3}
% —— 印面逐字符号路由补钉（探针实证 0914：书宋缺 U+2212，▱ 诸字库皆缺→NSC 子块；
%     禁改 sty 本体保 md5；无 ▱/− 用件的片可整块照抄，零副作用） ——
\xeCJKDeclareCharClass{Default}{"2212}%
\setCJKfamilyfont{nscblk}[Path=C:/提示词/工作区/字替对照-0909/variantF/fonts/,BoldFont=NSC-w600.ttf]{NSC-w500.ttf}%
\xeCJKDeclareSubCJKBlock{nscblk}{"25B1}%
\newcommand{\pxparallelogram}{{\CJKfamily{nscblk}▱}}%
% —— 断行微调（纯题档/详解档三零门：CJK 胶收缩 0.01→0.05em；应急伸展曾试 12em/2em，
%     2026-09-14 实测撤行回落 sty 默认 1em 三零仍守，故不设该行——补丁最少化；
%     只动伸缩分量不动自然宽度，版面纹理零漂；Underfull 治理读数见过程件） ——
\xeCJKsetup{CJKglue={\hskip 0.10em plus 0.08em minus 0.05em}}%
% —— 页脚件名（qp-headfoot 复用入口） ——
\renewcommand{\qpzhangming}{第二章\quad 平面解析几何}%
\renewcommand{\qpjianming}{导学件}%

\begin{document}

\zhangtitle{第二章\quad 平面解析几何}
\jietitle{2.1 坐标法}
\mubiaoline
\mubiaomu{1}{理解数轴上两点间距离与中点坐标公式，掌握平面内两点间距离公式与中点坐标公式，并能正确应用；}
\mubiaomu{2}{理解坐标法的基本思想，会用坐标法计算距离、中点等几何量，证明平行、垂直、对称等几何结论；}
\mubiaomu{3}{会用坐标法讨论存在性问题，体会“几何问题代数化、代数结果几何化”的解析几何方法价值.}

\vspace{1.7mm}
\begin{multicols}{2}
\emergencystretch=1em

% ============================================================
% 课前预习（知识点导学填空；零键零答案——预习键位按检查口令④裁定前口径，
% \kongbai 空线＋\zhentib 判断题面形，两档等形不泄答）
% ============================================================
\huaxing{课}{前}{预}{习}{知识导学\quad 素养初识}

\zsd{一}{数轴上的基本公式}

\tiaomu{1}{设\(A\)，\(B\)是数轴上坐标依次为\(x_1\)，\(x_2\)的两点，则有向线段\(\overrightarrow{AB}\)的数量\(AB=\)\kongbai{}；两点间距离\(d(A,B)=|x_1-x_2|\)；线段\(AB\)的中点坐标\(x_0=\)\kongbai{}．\zhuzhu{数量带符号、距离不带符号；中点公式是坐标和的一半，数轴对称点问题都归到中点式.}}

\zsd{二}{平面内两点的距离与中点}

\tiaomu{1}{设\(A(x_1,y_1)\)，\(B(x_2,y_2)\)，则两点间距离\(|AB|=\)\kongbai{}；线段\(AB\)的中点坐标为\kongbai{}．\zhuzhu{距离公式与中点公式是坐标法最基本的两个工具，本章各节都要反复使用；用前先核对两点坐标的符号.}}

\tiaomu{2}{由距离与中点反求坐标（或求参数）时，先设未知数，把几何条件写成\kongbai{}（方程），解出后回代检验．\zhuzhu{垂直与直角条件可经勾股定理（斜边平方等于两直角边平方和）转化为距离方程；存在性问题的根常不唯一，检验不可省.}}

\zsd{三}{坐标法}

\tiaomu{1}{坐标法（解析法）研究几何问题的基本思路：建立平面直角坐标系，设出点的坐标，把几何条件转化为\kongbai{}运算，算得结果后再回译成\kongbai{}结论．\zhuzhu{“几何代数化、代数几何化”是解析几何总纲；建系时选互相垂直的边所在直线、对称轴为坐标轴，可使运算最简.}}

\zhenhead{判断正误(正确的打\gou{},错误的打\cha{})}

\zhentib{(1)}{数轴上两点间的距离等于这两点坐标之和的绝对值．}

\zhentib{(2)}{平面内两点间的距离与按哪一点为起点计算有关．}

\liubai[9mm]{此处书写}

% ============================================================
% 课中探究（考点探究〔例+变式〕＝练/拓槽 18 键；题后紧跟答案＝ansblock 内嵌，
% 值照答案侧逐字，详解承定稿回嵌＋印面清洗 R1~R3）
% ============================================================
\huaxing{课}{中}{探}{究}{考点探究\quad 素养小结}

% ---------- 探究点一 距离与中点公式直用 ----------
\tjdnr{一}{距离与中点公式直用}{例\textbf{1}}{简单(知识点二)}{已知\(A(3,0)\)，\(B(0,-4)\)，求这两点之间的距离以及它们的中点坐标．}
\xiexwei{7mm}

\begin{ansblock}[2章-练-课时01-E4]
% ans:2章-练-课时01-E4
\ansitem{1}{|AB|＝5；中点坐标为(3/2,−2)．}
\ansnote{详解}{\(|AB|=\sqrt{(0-3)^{2}+(-4-0)^{2}}=\sqrt{9+16}=5\)；中点坐标为\(((3+0)/2,(0-4)/2)=(3/2,-2)\)．}
\end{ansblock}

\liB{变式\textbf{2}}{中档(知识点二)}{}{已知四边形\(MNPQ\)的顶点\(M(1,1)\)，\(N(3,-1)\)，\(P(4,0)\)，\(Q(2,2)\)，则四边形\(MNPQ\)的形状为\nobreak(\kongwei)}

\bindopt {\fontsize{10.5pt}{\qpTJgLeadLiti}\selectfont \makebox[\dimexpr0.5\linewidth-1em\relax][l]{A．平行四边形}\makebox[\dimexpr0.5\linewidth-1em\relax][l]{B．菱形}\par}

\bindopt {\fontsize{10.5pt}{\qpTJgLeadLiti}\selectfont \makebox[\dimexpr0.5\linewidth-1em\relax][l]{C．梯形}\makebox[\dimexpr0.5\linewidth-1em\relax][l]{D．矩形}\par}

\begin{ansblock}[2章-练-课时01-E1]
% ans:2章-练-课时01-E1
\ansitem{2}{D}
\ansnote{详解}{四边形\(MNPQ\)的对角线为\(MP\)与\(NQ\)．\(MP\)的中点为\(((1+4)/2,(1+0)/2)=(5/2,1/2)\)，\(NQ\)的中点为\(((3+2)/2,(-1+2)/2)=(5/2,1/2)\)，两对角线互相平分，故\(MNPQ\)是平行四边形．又 \(|MP|^{2}=(4-1)^{2}+(0-1)^{2}=10\)，\(|NQ|^{2}=(2-3)^{2}+(2+1)^{2}=10\)，对角线相等，故\pxparallelogram\(MNPQ\)为矩形．（若担心为正方形：\(|MN|^{2}=(3-1)^{2}+(-1-1)^{2}=8\)，\(|NP|^{2}=(4-3)^{2}+(0+1)^{2}=2\)，邻边不等，非正方形．）故选：D．}
\end{ansblock}

\liB{变式\textbf{3}}{简单(知识点一)}{}{已知数轴上\(A(-1)\)，\(B(2)\)，且\(A\)关于\(B\)的对称点为\(C\)，求\(C\)的坐标．}
\xiexwei{7mm}

\begin{ansblock}[2章-练-课时01-E5]
% ans:2章-练-课时01-E5
\ansitem{3}{C的坐标为5．}
\ansnote{详解}{\(B\)是线段\(AC\)的中点，设\(C(x)\)，则 \((-1+x)/2=2\)，解得\(x=5\)，即\(C\)的坐标为5．}
\end{ansblock}

\liB{变式\textbf{4}}{简单(知识点一)}{}{已知数轴上的点\(P\)到\(A(-9)\)的距离是它到\(B(-3)\)的距离的2倍，求点\(P\)的坐标．}
\xiexwei{7mm}

\begin{ansblock}[2章-练-课时01-E6]
% ans:2章-练-课时01-E6
\ansitem{4}{P的坐标为3或−5．}
\ansnote{详解}{设\(P(x)\)，依题意 \(|x+9|=2|x+3|\)．两边均非负，平方得 \((x+9)^{2}=4(x+3)^{2}\)，即 \(x^{2}+18x+81=4x^{2}+24x+36\)，整理得 \(3x^{2}+6x-45=0\)，即 \(x^{2}+2x-15=0\)，解得\(x=3\)或\(x=-5\)．检验：\(x=3\)时\(|3+9|=12=2|3+3|\)；\(x=-5\)时\(|-5+9|=4=2|-5+3|\)．所以点\(P\)的坐标为3或\(-5\)．}
\end{ansblock}

\tjdnr{一}{距离与中点公式直用}{例\textbf{5}}{中档(知识点二)}{已知\(\triangle ABC\)的三个顶点分别为\(A(2,1)\)，\(B(-2,3)\)，\(C(0,-1)\)，求\(\triangle ABC\)三条中线的长．}
\xiexwei{12mm}

\begin{ansblock}[2章-练-课时01-E9]
% ans:2章-练-课时01-E9
\ansitem{5}{AB边上的中线长3；BC边上的中线长3；CA边上的中线长3√2．}
\ansnote{详解}{设\(AB\)、\(BC\)、\(CA\)的中点分别为\(M\)、\(N\)、\(P\)．\(M((2-2)/2,(1+3)/2)=M(0,2)\)，中线\(|CM|=\sqrt{(0-0)^{2}+(2+1)^{2}}=3\)；\(N((-2+0)/2,(3-1)/2)=N(-1,1)\)，中线\(|AN|=\sqrt{(-1-2)^{2}+(1-1)^{2}}=3\)；\(P((0+2)/2,(-1+1)/2)=P(1,0)\)，中线\(|BP|=\sqrt{(1+2)^{2}+(0-3)^{2}}=\sqrt{18}=3\sqrt{2}\)．所以三条中线的长分别为3、3、\(3\sqrt{2}\)．}
\end{ansblock}

\liB{变式\textbf{6}}{中档(知识点二)}{}{已知\(A(-1,3)\)，\(B(3,3)\)，\(C\left(1,2\sqrt{3}+3\right)\)，证明\(\triangle ABC\)是等边三角形．}
\xiexwei{12mm}

\begin{ansblock}[2章-练-课时01-E8]
% ans:2章-练-课时01-E8
\ansitem{6}{证明见详解．}
\ansnote{详解}{由两点间距离公式：\par\noindent\(|AB|=\sqrt{(3+1)^{2}+(3-3)^{2}}=\sqrt{16}=4\)；\par\noindent\(|AC|=\sqrt{(1+1)^{2}+(2\sqrt{3}+3-3)^{2}}=\sqrt{4+12}=4\)；\par\noindent\(|BC|=\sqrt{(1-3)^{2}+(2\sqrt{3}+3-3)^{2}}=\sqrt{4+12}=4\)．\par\noindent 所以\(|AB|=|AC|=|BC|\)，故\(\triangle ABC\)是等边三角形．}
\end{ansblock}

\tjdnr{一}{距离与中点公式直用}{例\textbf{7}}{中档(知识点二)}{定义点到曲线的距离为该点与曲线上所有点之间距离的最小值，则点\(P(0,1)\)到曲线 \(x-xy+2=0\) 的距离为\kongbai{}．}
\xiexwei{7mm}

\begin{ansblock}[2章-练-课时01-E7]
% ans:2章-练-课时01-E7
\ansitem{7}{2}
\ansnote{详解}{曲线方程中\(x=0\)时\(2=0\)不成立，故\(x\neq 0\)，且 \(y=1+2/x\)．设曲线上任意一点\(A(x,1+2/x)\)，则\(|AP|=\sqrt{x^{2}+(1+2/x-1)^{2}}=\sqrt{x^{2}+4/x^{2}}\)．由基本不等式 \(x^{2}+4/x^{2}\geqslant 2\sqrt{x^{2}\cdot 4/x^{2}}=4\)，当且仅当 \(x^{2}=4/x^{2}\)，即\(x=\pm\sqrt{2}\) 时等号成立，此时\(|AP|\)最小值为\(\sqrt{4}=2\)．故所求距离为2．}
\end{ansblock}

\liB{变式\textbf{8}}{中档(知识点二)}{\duoxuan\hspace{0.6em}}{定义点\(P(x_0,y_0)\)到直线\(l\)：\(ax+by+c=0\)（\(a^{2}+b^{2}\neq 0\)）的有向距离为 \(d=(ax_0+by_0+c)/\sqrt{a^{2}+b^{2}}\)．已知点\(P_1\)，\(P_2\)到直线\(l\)的有向距离分别是\(d_1\)，\(d_2\)，则以下命题不正确的是\nobreak(\kongwei)}

\lxopt{A．若\(d_1=d_2=1\)，则直线\(P_1P_2\)与直线\(l\)平行}

\lxopt{B．若\(d_1=1\)，\(d_2=-1\)，则直线\(P_1P_2\)与直线\(l\)垂直}

\lxopt{C．若\(d_1+d_2=0\)，则直线\(P_1P_2\)与直线\(l\)垂直}

\lxopt{D．若\(d_1\cdot d_2\leqslant 0\)，则直线\(P_1P_2\)与直线\(l\)相交}

{\ansblockgrayfalse
\begin{ansblock}[2章-练-课时01-E2]
% ans:2章-练-课时01-E2
\ansitem{8}{BCD}
\ansnote{详解}{设\(P_1(x_1,y_1)\)，\(P_2(x_2,y_2)\)．A：\(d_1=d_2=1\)，则 \(ax_1+by_1+c=ax_2+by_2+c=\sqrt{a^{2}+b^{2}}\)，即 \(P_1\)、\(P_2\) 都在直线 \(ax+by+c-\sqrt{a^{2}+b^{2}}=0\) 上，该直线与\(l\)平行，故直线\(P_1P_2\parallel l\)，A正确；B：\(d_1=1\)，\(d_2=-1\) 只说明\(P_1\)、\(P_2\)在\(l\)的两侧且到\(l\)的距离相等，线段\(P_1P_2\)的中点落在\(l\)上，但\(P_1P_2\)与\(l\)不一定垂直，B错误；C：\(d_1+d_2=0\) 包含\(d_1=d_2=0\)的情形，此时\(P_1\)、\(P_2\)都在\(l\)上，直线\(P_1P_2\)与\(l\)重合，谈不上垂直，C错误；D：\(d_1\cdot d_2\leqslant 0\) 表示\(P_1\)、\(P_2\)分居\(l\)两侧或至少一点在\(l\)上，此时直线\(P_1P_2\)与\(l\)相交或重合，D把“重合”漏掉，D错误．故选：BCD．}
\end{ansblock}}

\liB{变式\textbf{9}}{中档(知识点二)}{\duoxuan\hspace{0.6em}}{在平面直角坐标系中，已知点\(A(x_1,y_1)\)、\(B(x_2,y_2)\)，定义 \(d(A,B)=|x_1-x_2|+|y_1-y_2|\) 为两点\(A\)，\(B\)的“折线距离”；又设点\(P\)及直线\(l\)上任意一点\(Q\)，称\(d(P,Q)\)的最小值为点\(P\)到直线\(l\)的“折线距离”，记作\(d(P,l)\)．则下列说法正确的是\nobreak(\kongwei)}

\lxopt{A．对任意的两点\(A\)，\(B\)，都有 \(d(A,B)\geqslant|AB|\)}

\lxopt{B．对任意三点\(A\)、\(B\)、\(C\)，都有 \(d(A,C)+d(B,C)\geqslant d(A,B)\)}

\lxopt{C．已知点\(P(3,1)\)和直线\(l\)：\(2x-y-1=0\)，则 \(d(P,l)=4\sqrt{5}/5\)}

\lxopt{D．已知点\(O(0,0)\)，动点\(P(x,y)\)满足\(d(P,O)=1\)，则动点\(P\)的轨迹围成平面图形的面积是2}

{\ansblockgrayfalse
\begin{ansblock}[2章-练-课时01-E3]
% ans:2章-练-课时01-E3
\ansitem{9}{ABD}
\ansnote{详解}{A：\((|x_1-x_2|+|y_1-y_2|)^{2}=(x_1-x_2)^{2}+2|x_1-x_2||y_1-y_2|+(y_1-y_2)^{2}\geqslant(x_1-x_2)^{2}+(y_1-y_2)^{2}\)，开方即得 \(d(A,B)\geqslant|AB|\)，A正确；B：由绝对值三角不等式，\(|x_1-x_2|+|x_2-x_3|\geqslant|x_1-x_3|\)，\(|y_1-y_2|+|y_2-y_3|\geqslant|y_1-y_3|\)，两式相加即得 \(d(A,B)+d(B,C)\geqslant d(A,C)\)（将\(A\)、\(B\)、\(C\)重排即题设形式），B正确；C：设\(Q(x,2x-1)\)是\(l\)上任一点，\(d(P,Q)=|x-3|+|2x-2|=|x-3|+|x-1|+|x-1|\geqslant|(x-3)-(x-1)|+0=2\)，当且仅当\(x=1\)时等号成立，故 \(d(P,l)=2\neq 4\sqrt{5}/5\)，C错误；D：\(d(P,O)=|x|+|y|=1\) 的图象是以\((1,0)\)、\((0,1)\)、\((-1,0)\)、\((0,-1)\)为顶点的正方形，对角线长均为2，面积\(=(1/2)\times 2\times 2=2\)，D正确．故选：ABD．}
\end{ansblock}}

\xiaojie{①识别：以求两点间距离、中点坐标，或由距离/中点条件反求坐标（参数）、判定图形形状为目标的题，判归本题型；新定义距离类先把定义式翻译成坐标运算再套用.}
\bindp {\kaishu ②操作：建系设点，直写两点间距离公式与中点坐标公式；数轴问题用绝对值距离与一次式中点式，平方去绝对值后须回代检验.}
\bindp {\kaishu ③收束：判定四边形形状先看对角线——互相平分定平行四边形，再证相等定矩形、垂直定菱形；反求坐标的根不唯一时逐根检验.}

% ---------- 探究点二 坐标法证明与存在性讨论 ----------
\tjdnr{二}{坐标法证明与存在性讨论}{例\textbf{10}}{中档(知识点三)}{已知\(ABCD\)是一个长方形，且\(M\)是\(ABCD\)所在平面上任意一个点，求证：\(|AM|^{2}+|CM|^{2}=|BM|^{2}+|DM|^{2}\)．}
\xiexwei{14mm}

\begin{ansblock}[2章-练-课时01-E11]
% ans:2章-练-课时01-E11
\ansitem{10}{证明见详解．}
\ansnote{详解}{以长方形的顶点\(A\)为原点、\(AB\)所在直线为\(x\)轴建立平面直角坐标系．设\(AB=a\)，\(AD=b\)，则\(A(0,0)\)，\(B(a,0)\)，\(C(a,b)\)，\(D(0,b)\)．再设\(M(x,y)\)．\(|AM|^{2}+|CM|^{2}=x^{2}+y^{2}+(x-a)^{2}+(y-b)^{2}\)；\(|BM|^{2}+|DM|^{2}=(x-a)^{2}+y^{2}+x^{2}+(y-b)^{2}\)．两式展开后均为 \(x^{2}+y^{2}+(x-a)^{2}+(y-b)^{2}\)，恒相等．所以 \(|AM|^{2}+|CM|^{2}=|BM|^{2}+|DM|^{2}\)，即对平面上任意一点\(M\)结论都成立．}
\end{ansblock}

\liB{变式\textbf{11}}{中档(知识点三)}{}{求证：平行四边形两条对角线的平方和等于四条边的平方和．}
\xiexwei{14mm}

{\ansblockgrayfalse
\begin{ansblock}[2章-练-课时01-E14]
% ans:2章-练-课时01-E14
\ansitem{11}{证明见详解．}
\ansnote{详解}{以平行四边形相邻两边为坐标轴方向建立平面直角坐标系：取顶点\(A\)为原点、\(AB\)所在直线为\(x\)轴．设\(B(a,0)\)，\(D(b,c)\)，则由平行四边形性质知 \(C(a+b,c)\)．两条对角线：\(|AC|^{2}=(a+b)^{2}+c^{2}\)，\(|BD|^{2}=(b-a)^{2}+c^{2}\)，所以 \(|AC|^{2}+|BD|^{2}=\left(a^{2}+2ab+b^{2}+c^{2}\right)+\left(a^{2}-2ab+b^{2}+c^{2}\right)=2a^{2}+2b^{2}+2c^{2}\)．四条边：\(|AB|^{2}=a^{2}\)，\(|CD|^{2}=a^{2}\)，\(|AD|^{2}=b^{2}+c^{2}\)，\(|BC|^{2}=b^{2}+c^{2}\)，所以 \(|AB|^{2}+|BC|^{2}+|CD|^{2}+|DA|^{2}=2a^{2}+2b^{2}+2c^{2}\)．两式相同，故平行四边形两条对角线的平方和等于四条边的平方和．}
\end{ansblock}}

\tjdnr{二}{坐标法证明与存在性讨论}{例\textbf{12}}{中档(知识点三)}{已知正方形\(ABCD\)的边长为4，若\(E\)是\(BC\)的中点，\(F\)是\(CD\)的中点，求证：\(BF\perp AE\)．}
\xiexwei{14mm}

\begin{ansblock}[2章-练-课时01-E13]
% ans:2章-练-课时01-E13
\ansitem{12}{证明见详解．}
\ansnote{详解}{以\(A\)为原点、\(AB\)所在直线为\(x\)轴建立平面直角坐标系，则\(A(0,0)\)，\(B(4,0)\)，\(C(4,4)\)，\(D(0,4)\)；由\(E\)、\(F\)分别是\(BC\)、\(CD\)的中点得 \(E(4,2)\)，\(F(2,4)\)．于是 \(\overrightarrow{AE}=(4,2)\)，\(\overrightarrow{BF}=(2-4,4-0)=(-2,4)\)．\(\overrightarrow{AE}\cdot\overrightarrow{BF}=4\times(-2)+2\times 4=-8+8=0\)，所以 \(\overrightarrow{AE}\perp\overrightarrow{BF}\)，即 \(BF\perp AE\)．}
\end{ansblock}

\liB{变式\textbf{13}}{中档(知识点三)}{}{已知\(A(3,1)\)，\(B(-2,2)\)，在\(y\)轴上的点\(P\)满足\(PA\perp PB\)，求\(P\)的坐标．}
\xiexwei{12mm}

\begin{ansblock}[2章-练-课时01-E10]
% ans:2章-练-课时01-E10
\ansitem{13}{P(0,4)或P(0,−1)．}
\ansnote{详解}{设\(P(0,y)\)．\(PA\perp PB\) 即\(\triangle PAB\)是以\(AB\)为斜边的直角三角形，\(|PA|^{2}+|PB|^{2}=|AB|^{2}\)．\(|PA|^{2}=9+(y-1)^{2}\)，\(|PB|^{2}=4+(y-2)^{2}\)，\(|AB|^{2}=(-2-3)^{2}+(2-1)^{2}=26\)，代入：\(9+y^{2}-2y+1+4+y^{2}-4y+4=26\)，整理得 \(2y^{2}-6y-8=0\)，即 \(y^{2}-3y-4=0\)，解得\(y=4\)或\(y=-1\)．所以点\(P\)的坐标为\((0,4)\)或\((0,-1)\)．}
\end{ansblock}

\tjdnr{二}{坐标法证明与存在性讨论}{例\textbf{14}}{中档(知识点三)}{已知函数\(y=f(x)\)的图象与函数\(y=x^{2}+1\)的图象关于点\(M(2,0)\)对称，求\(f(x)\)的解析式．}
\xiexwei{12mm}

\begin{ansblock}[2章-练-课时01-E12]
% ans:2章-练-课时01-E12
\ansitem{14}{f(x)＝−(x−4)²−1．}
\ansnote{详解}{在\(y=f(x)\)的图象上任取一点\(P(x,y)\)，它关于点\(M(2,0)\)的对称点为 \(P'(4-x,-y)\)．由对称性，\(P'\)在函数\(y=x^{2}+1\)的图象上，所以 \(-y=(4-x)^{2}+1\)，即 \(y=-(x-4)^{2}-1\)．故 \(f(x)=-(x-4)^{2}-1\)．}
\end{ansblock}

\tjdnr{二}{坐标法证明与存在性讨论}{例\textbf{15}}{简单(知识点三)}{在 Rt\(\triangle ABC\) 中，\(\angle C=90^\circ\)，以 \(C\) 为原点、\(CB\)、\(CA\) 所在直线分别为 \(x\) 轴、\(y\) 轴建立平面直角坐标系，设 \(B(b,0)\)、\(A(0,a)\)（\(a>0\)，\(b>0\)），\(M\) 为斜边 \(AB\) 的中点，则下列结论正确的是\nobreak(\kongwei)}

\bindopt {\fontsize{10.5pt}{\qpTJgLeadLiti}\selectfont \makebox[\dimexpr0.5\linewidth-1em\relax][l]{A．\(|CM|=|AB|\)}\makebox[\dimexpr0.5\linewidth-1em\relax][l]{B．\(|CM|=|AB|/2\)}\par}

\bindopt {\fontsize{10.5pt}{\qpTJgLeadLiti}\selectfont \makebox[\dimexpr0.5\linewidth-1em\relax][l]{C．\(|CM|=(a+b)/2\)}\makebox[\dimexpr0.5\linewidth-1em\relax][l]{D．\(|CM|=(a^{2}+b^{2})/2\)}\par}

\begin{ansblock}[2章-练-课时01-E15]
% ans:2章-练-课时01-E15
\ansitem{15}{B}
\ansnote{详解}{由中点坐标公式，\(M((b+0)/2,(0+a)/2)\)，即 \(M(b/2,a/2)\)。由两点间距离公式：\(|CM|=\sqrt{(b/2)^{2}+(a/2)^{2}}=\sqrt{a^{2}+b^{2}}/2\)；\(|AB|=\sqrt{(b-0)^{2}+(0-a)^{2}}=\sqrt{a^{2}+b^{2}}\)。故 \(|CM|=|AB|/2\)，选 B。这就用坐标法证明了：直角三角形斜边上的中线等于斜边的一半。}
\end{ansblock}

\liB{变式\textbf{16}}{简单(知识点三)}{}{在\(\triangle ABC\) 中，以 \(B\) 为原点、\(BC\) 所在直线为 \(x\) 轴建立平面直角坐标系，设 \(C(a,0)\)、\(A(b,c)\)（\(a>0\)，\(c>0\)）。\(M\)、\(N\) 分别为 \(AB\)、\(AC\) 的中点，则下列关系成立的是\nobreak(\kongwei)}

\lxopt{A．\(MN\parallel BC\)，且 \(|MN|=|BC|\)}

\lxopt{B．\(MN\parallel BC\)，且 \(|MN|=|BC|/2\)}

\lxopt{C．\(MN\perp BC\)，且 \(|MN|=|BC|/2\)}

\lxopt{D．\(MN\) 与 \(BC\) 不平行，且 \(|MN|=|BC|/2\)}

\begin{ansblock}[2章-练-课时01-E16]
% ans:2章-练-课时01-E16
\ansitem{16}{B}
\ansnote{详解}{由中点坐标公式，\(M(b/2,c/2)\)，\(N((a+b)/2,c/2)\)。两点纵坐标相等且 \(|MN|=a/2>0\)（\(M\neq N\)），故 \(MN\parallel x\) 轴；而 \(BC\) 在 \(x\) 轴上，所以 \(MN\parallel BC\)。又 \(|MN|=|(a+b)/2-b/2|=a/2\)，\(|BC|=a\)，故 \(|MN|=|BC|/2\)。选 B。这就用坐标法证明了：三角形中位线平行于第三边，且等于第三边的一半。}
\end{ansblock}

\xiaojie{①识别：凡“求证”类几何结论与“是否存在/能否”类讨论，判归本题型.}
\bindp {\kaishu ②操作：选合适坐标系（以互相垂直的边所在直线为轴、对称中心为原点），设点坐标，把几何条件译成坐标等式（距离、向量数量积、对称关系）后再化简.}
\bindp {\kaishu ③收束：存在性问题先设参再解方程，解出后必须回代检验；证明题最后一句要回译成几何结论.}

% ---------- 探究点三 拓展·收难 ----------
\tjdnr{三}{拓展·收难}{例\textbf{17}}{难(知识点三)}{在数轴上，对坐标分别为\(x_1\)和\(x_2\)的两点\(A\)和\(B\)，用绝对值定义两点间的距离，表示为 \(d(A,B)=|x_1-x_2|\)．}

\bindp (1) 在数轴上任意取三点\(A\)，\(B\)，\(C\)，证明 \(d(A,B)\leqslant d(A,C)+d(B,C)\)；

\bindp (2) 设\(A\)和\(B\)两点的坐标分别为\(-3\)和\(2\)，分别找出(1)中不等式等号成立和等号不成立时点\(C\)的范围．
\xiexwei{18mm}

{\ansblockgrayfalse
\begin{ansblock}[2章-拓-课时01-T1]
% ans:2章-拓-课时01-T1
\ansitem{17}{(1)见详解；(2) 等号成立$\iff$C的坐标x∈[−3,2]；等号不成立$\iff$x<−3或x>2．}
\ansnote{详解}{(1) 设\(C\)的坐标为\(x_3\)．因为 \(x_1-x_2=(x_1-x_3)+(x_3-x_2)\)，由绝对值三角不等式\(d(A,B)=|x_1-x_2|=|(x_1-x_3)+(x_3-x_2)|\leqslant|x_1-x_3|+|x_3-x_2|=d(A,C)+d(B,C)\)，得证．(2) \(|x_1-x_2|=|(x_1-x_3)+(x_3-x_2)|\)取等号\(\iff(x_1-x_3)(x_3-x_2)\geqslant 0\)，即\(x_3\)介于\(x_1\)与\(x_2\)之间（含端点）．此处\(x_1=-3\)，\(x_2=2\)，故等号成立\(\iff -3\leqslant x\leqslant 2\)；等号不成立（即\(d(A,B)<d(A,C)+d(B,C)\)）\(\iff x<-3\)或\(x>2\)．}
\end{ansblock}}

\liB{变式\textbf{18}}{难(知识点三)}{}{城市的许多街道是相互垂直或平行的，往往不能沿直线行走到达目的地，只能按直角拐弯的方式行走．按照街道的垂直和平行方向建立平面直角坐标系，对两点\(A(x_1,y_1)\)和\(B(x_2,y_2)\)，定义两点间距离为 \(d(A,B)=|x_1-x_2|+|y_1-y_2|\)．}

\bindp (1) 在平面直角坐标系中任意取三点\(A\)，\(B\)，\(C\)，证明 \(d(A,B)\leqslant d(A,C)+d(B,C)\)；

\bindp (2) 设\(A(x_1,y_1)\)，\(B(x_2,y_2)\)，分别找出(1)中不等式等号成立和等号不成立时点\(C\)的范围．
\xiexwei{18mm}

{\ansblockgrayfalse
\begin{ansblock}[2章-拓-课时01-T2]
% ans:2章-拓-课时01-T2
\ansitem{18}{(1)见详解；(2) 等号成立$\iff$C在以AB为对角线的矩形内（含边界），即$(x_C-x_1)(x_C-x_2)\leqslant 0$且$(y_C-y_1)(y_C-y_2)\leqslant 0$；等号不成立$\iff$C在该矩形外．}
\ansnote{详解}{(1) 设\(C(x_3,y_3)\)．\(x_1-x_2=(x_1-x_3)+(x_3-x_2)\)，\(y_1-y_2=(y_1-y_3)+(y_3-y_2)\)，两次用绝对值三角不等式：\(d(A,B)=|x_1-x_2|+|y_1-y_2|\leqslant(|x_1-x_3|+|x_3-x_2|)+(|y_1-y_3|+|y_3-y_2|)=d(A,C)+d(B,C)\)，得证．(2) 横、纵两个方向分别取等：第一个不等式取等\(\iff(x_1-x_3)(x_3-x_2)\geqslant 0\)，第二个取等\(\iff(y_1-y_3)(y_3-y_2)\geqslant 0\)，故总等号成立\(\iff x_3\)介于\(x_1\)、\(x_2\)之间且\(y_3\)介于\(y_1\)、\(y_2\)之间，即\(C\)落在以\(AB\)为对角线（边平行于坐标轴）的矩形内部或边界上；只要有一个方向不取等，就有\(d(A,B)<d(A,C)+d(B,C)\)，即\(C\)在该矩形外．}
\end{ansblock}}

\xiaojie{①识别：以绝对值距离、折线距离等新定义出发的不等式证明与取等范围讨论，判归本题型.}
\bindp {\kaishu ②操作：核心工具是绝对值三角不等式\(|a+b|\leqslant|a|+|b|\)；把差拆成两段之和，两次放缩即得证明.}
\bindp {\kaishu ③收束：取等当且仅当两段同号或其一为零；多维情形逐分量独立取等，范围结论按分量乘积（区域）表述.}

% ============================================================
% 课堂评价 G5（恒5＝3单选+2填空＝导槽 G1~G5；ketang 新制顺排可跨栏，禁 \vbox 装栏）
% 印面号 19—23；多选门：本片正文多选 2（E2/E3）≤4 合规
% ============================================================
\begin{ketang}{本组共5题（第19—23题），19—21为单选，22—23为填空．}

\jiancestem{{\fontsize{11.4pt}{14pt}\selectfont\heihao {\numboldjian 19}．}\hspace{4.1mm}\tieside{简单(知识点二)}\hspace{2.2mm}点\(P(-2,5)\)为平面直角坐标系内一点，线段\(PM\)的中点是\((1,0)\)，那么点\(M\)到原点\(O\)的距离为\nobreak(\kongwei)}

\bindopt {\fontsize{10.5pt}{\qpTJgLeadPingjia}\selectfont \makebox[\dimexpr0.25\linewidth-0.5em\relax][l]{A．\(41\)}\makebox[\dimexpr0.25\linewidth-0.5em\relax][l]{B．\(\sqrt{41}\)}\makebox[\dimexpr0.25\linewidth-0.5em\relax][l]{C．\(\sqrt{39}\)}\makebox[\dimexpr0.25\linewidth-0.5em\relax][l]{D．\(39\)}\par}

\begin{ansblock}[2章-导-课时01-G1]
% ans:2章-导-课时01-G1
\ansitem{19}{B}
\ansnote{详解}{设\(M(x,y)\)，由中点坐标公式得 \((x-2)/2=1\)，\((y+5)/2=0\)，解得\(x=4\)，\(y=-5\)，所以点\(M(4,-5)\)．则 \(|OM|=\sqrt{4^{2}+(-5)^{2}}=\sqrt{41}\)．故选：B．}
\end{ansblock}

\jiancestem{{\fontsize{11.4pt}{14pt}\selectfont\heihao {\numboldjian 20}．}\hspace{4.1mm}\tieside{简单(知识点二)}\hspace{2.2mm}以\(A(-1,1)\)，\(B(2,-1)\)，\(C(1,4)\)为顶点的三角形是\nobreak(\kongwei)}

\lxopt{A．以\(A\)点为直角顶点的直角三角形}

\lxopt{B．以\(B\)点为直角顶点的直角三角形}

\lxopt{C．锐角三角形}

\lxopt{D．钝角三角形}

\begin{ansblock}[2章-导-课时01-G2]
% ans:2章-导-课时01-G2
\ansitem{20}{A}
\ansnote{详解}{由两点间距离公式：\(|AB|^{2}=(2+1)^{2}+(-1-1)^{2}=9+4=13\)；\(|AC|^{2}=(1+1)^{2}+(4-1)^{2}=4+9=13\)；\(|BC|^{2}=(1-2)^{2}+(4+1)^{2}=1+25=26\)．因为\(|AB|^{2}+|AC|^{2}=13+13=26=|BC|^{2}\)，由勾股定理逆定理知\(\triangle ABC\)是以\(A\)为直角顶点的直角三角形（且\(|AB|=|AC|\)，还是等腰直角三角形）．故选：A．}
\end{ansblock}

\jiancestem{{\fontsize{11.4pt}{14pt}\selectfont\heihao {\numboldjian 21}．}\hspace{4.1mm}\tieside{中档(知识点二)}\hspace{2.2mm}在平面直角坐标系内有两个点\(A(4,2)\)，\(B(1,-2)\)，若在\(x\)轴上存在点\(C\)，使\(\angle ACB=\pi/2\)，则点\(C\)的坐标是\nobreak(\kongwei)}

\bindopt {\fontsize{10.5pt}{\qpTJgLeadPingjia}\selectfont \makebox[\dimexpr0.5\linewidth-1em\relax][l]{A．\((3,0)\)}\makebox[\dimexpr0.5\linewidth-1em\relax][l]{B．\((0,0)\)}\par}

\bindopt {\fontsize{10.5pt}{\qpTJgLeadPingjia}\selectfont \makebox[\dimexpr0.5\linewidth-1em\relax][l]{C．\((5,0)\)}\makebox[\dimexpr0.5\linewidth-1em\relax][l]{D．\((0,0)\)或\((5,0)\)}\par}

\begin{ansblock}[2章-导-课时01-G3]
% ans:2章-导-课时01-G3
\ansitem{21}{D}
\ansnote{详解}{设\(C(x_0,0)\)．\(\angle ACB=\pi/2\) 即\(\triangle ACB\)是以\(AB\)为斜边的直角三角形，故 \(|CA|^{2}+|CB|^{2}=|AB|^{2}\)．\(|CA|^{2}=(x_0-4)^{2}+4\)，\(|CB|^{2}=(x_0-1)^{2}+4\)，\(|AB|^{2}=(4-1)^{2}+(2+2)^{2}=9+16=25\)，代入得 \((x_0-4)^{2}+(x_0-1)^{2}+8=25\)，展开整理得 \(x_0^{2}-5x_0=0\)，解得\(x_0=0\)或\(x_0=5\)，经检验均使\(A\)、\(B\)、\(C\)构成三角形．所以点\(C\)的坐标为\((0,0)\)或\((5,0)\)．故选：D．}
\end{ansblock}

\jiancestem{{\fontsize{11.4pt}{14pt}\selectfont\heihao {\numboldjian 22}．}\hspace{4.1mm}\tieside{简单(知识点一)}\hspace{2.2mm}已知数轴上\(A(-2)\)，\(B(10)\)，求这两点之间的距离以及它们的中点坐标．}
\xiexwei{7mm}

\begin{ansblock}[2章-导-课时01-G4]
% ans:2章-导-课时01-G4
\ansitem{22}{距离为12；中点坐标为4．}
\ansnote{详解}{数轴上两点距离 \(d(A,B)=|10-(-2)|=12\)；中点坐标为 \((-2+10)/2=4\)．}
\end{ansblock}

\jiancestem{{\fontsize{11.4pt}{14pt}\selectfont\heihao {\numboldjian 23}．}\hspace{4.1mm}\tieside{简单(知识点二)}\hspace{2.2mm}已知\(A(a,0)\)，\(B(0,10)\)，且\(|AB|=17\)，求\(a\)的值．}
\xiexwei{7mm}

\begin{ansblock}[2章-导-课时01-G5]
% ans:2章-导-课时01-G5
\ansitem{23}{a＝±3√21．}
\ansnote{详解}{由两点间距离公式 \(|AB|=\sqrt{a^{2}+10^{2}}=17\)，故 \(a^{2}+100=289\)，\(a^{2}=189\)，\(a=\pm\sqrt{189}=\pm 3\sqrt{21}\)．}
\end{ansblock}

\end{ketang}

% ---- 尾块（\tailfill 置 \end{multicols} 前；无参＝内容尾随框，件侧不擅自定高） ----
\tailfill

\end{multicols}
\end{document}
